% TOP_PROBLEM: {"id": 26, "title": "Are there infinitely many Mersenne primes?", "category_id": 1, "category": {"id": 1, "name": "number_theory", "display_name": "Number Theory", "description": "Properties of integers, prime numbers, Diophantine equations.", "slug": "number-theory", "order_index": 1, "created_at": "2026-07-31T15:26:25.670Z"}, "status": "open"}
% FIRST_POSTED: 2026-09-25
% ENTRY_KIND: statement_only
% !TeX program = lualatex
% Definition and sources only; this file is not an LLM solution attempt.
\documentclass[11pt]{article}
\usepackage[margin=25mm]{geometry}
\usepackage{fontspec}
\setmainfont{Latin Modern Roman}
\usepackage{amsmath,amssymb,mathtools}
\usepackage{unicode-math}
\setmathfont{Latin Modern Math}
\usepackage{xurl,hyperref}
\hypersetup{colorlinks=true,urlcolor=blue}
\setcounter{secnumdepth}{0}
\setlength{\parindent}{0pt}
\setlength{\parskip}{5pt}
\setlength{\emergencystretch}{3em}
\begin{document}
\section*{\texorpdfstring{Are there infinitely many Mersenne primes?}{Are there infinitely many Mersenne primes?}}\label{26}
\subsection{Definitions and mathematical statement}
A Mersenne number is $M_p=2^p-1$. The selected infinitude question is
\[
 \forall B>0\quad\exists p>B\quad
 p\text{ prime and }2^p-1\text{ prime}.
\]
Primality of the exponent is necessary, but not sufficient, for a Mersenne number to be prime. The catalogue's equivalent even-perfect-number formulation uses the Euclid--Euler correspondence
$N=2^{p-1}(2^p-1)$ between Mersenne primes and even perfect numbers. The question requires neither a density prediction nor an efficient procedure guaranteed to find the next example. Arbitrarily large verified examples alone do not prove infinitude without a general argument.
\par\smallskip\textit{Formulation and concept references:} \href{https://t5k.org/mersenne/index.html}{[S1]}.

\subsection{Short English statement}
Are there infinitely many primes of the form $2^p-1$ with $p$ prime, equivalently infinitely many even perfect numbers?
\subsection{Sources}
\begin{itemize}
\item[S1]"Mersenne Primes: History, Theorems and Lists," The PrimePages (t5k.org), maintained curated reference; printed from the PrimePages, (c) Reginald McLean.
\url{https://t5k.org/mersenne/index.html}.
\textit{Formulation source.}
\end{itemize}
\end{document}
